% GENERATED FILE. Edit canonical lessons, then run npm run generate:books.
% canonical-source-hash: fnv1a64:1e1b07bdf895b685
% canonical-lessons: ML-C507-mansun, ML-C507-kalavarsam, ML-C507-tulavarsam, ML-C507-toruka, ML-C507-carralmala

\chapter{The Monsoon}
\label{ch:ml-the-monsoon}
\hlchaptermodality{507}

\begin{chapteropening}
\textbf{By the end of this chapter:} \emph{I can say the monsoon, the southwest monsoon, the northeast monsoon, to stop raining and a light drizzle.}

Complete the last lesson of chapter 507: I can say the monsoon, the southwest monsoon, the northeast monsoon, to stop raining and a light drizzle.
\end{chapteropening}

\section[\texorpdfstring{\ml{മൺസൂൺ} (maṇsūṇ)}{മൺസൂൺ (maṇsūṇ)}]{\ml{മൺസൂൺ} (maṇsūṇ) — the monsoon}

\label{lesson:ML-C507-mansun}



\begin{quote}
Before the new one: say the Malayalam for the atmosphere, then the Malayalam for rain clouds.
\end{quote}

\subsection*{You'll want to know: \ml{മൺസൂൺ}}
\textbf{\ml{മൺസൂൺ}} — \emph{maṇsūṇ} — \textquotedblleft{}the monsoon\textquotedblright{}.

\begin{quote}
\ml{ജൂണിൽ} \ml{മൺസൂൺ} \ml{തുടങ്ങും}.
\end{quote}

\begin{quote}
\emph{jūṇil maṇsūṇ tuṭaṅṅuṁ.}
\end{quote}

\begin{quote}
The monsoon will begin in June.
\end{quote}

English uses it for the season of the rains; Malayalam also says \ml{മഴക്കാലം}.

\begin{cousinweb}
\textbf{\ml{മൺസൂൺ}} (\emph{maṇsūṇ}) was borrowed from English \emph{monsoon}, the seasonal wind and the rains it brings.
\end{cousinweb}

\begin{grammarlens}[title={What you've built}]
One more word for getting around.
\end{grammarlens}

\subsection*{Your turn}
\begin{itemize}
\raggedright
  \item \emph{Say it:} \emph{maṇsūṇ}
  \item \emph{Say it:} \emph{maṇsūṇ}, once more
  \item \emph{Say it:} \emph{jūṇil maṇsūṇ tuṭaṅṅuṁ.} — The monsoon will begin in June.
  \item \emph{Recall:} say the Malayalam for the atmosphere, then the Malayalam for rain clouds, then say \emph{maṇsūṇ} again
\end{itemize}

\begin{tcolorbox}[breakable,colback=teal!4,colframe=teal!35!black,title={Before you move on}]
What does \ml{മൺസൂൺ} mean? (\textquotedblleft{}The monsoon\textquotedblright{}.) Say it once more.
\end{tcolorbox}
\section[\texorpdfstring{\ml{കാലവർഷം} (kālavarṣaṁ)}{കാലവർഷം (kālavarṣaṁ)}]{\ml{കാലവർഷം} (kālavarṣaṁ) — the southwest monsoon}

\label{lesson:ML-C507-kalavarsam}



\begin{quote}
Before the new one: say the Malayalam for rain clouds, then the Malayalam for the monsoon.
\end{quote}

\subsection*{You'll want to know: \ml{കാലവർഷം}}
\textbf{\ml{കാലവർഷം}} — \emph{kālavarṣaṁ} — \textquotedblleft{}the southwest monsoon\textquotedblright{}.

\begin{quote}
\ml{കാലവർഷം} \ml{ആ} \ml{വർഷം} \ml{നേരത്തെ} \ml{എത്തി}.
\end{quote}

\begin{quote}
\emph{kālavarṣaṁ ā varṣaṁ nēratte etti.}
\end{quote}

\begin{quote}
The southwest monsoon came early that year.
\end{quote}

The big rains of June and July, the first of the two monsoons.

\begin{cousinweb}
\textbf{\ml{കാലവർഷം}} (\emph{kālavarṣaṁ}) is a Malayalam compound built from two Sanskrit words, \emph{kāla}, \textquotedblleft{}time, season\textquotedblright{}, and \emph{varṣa}, \textquotedblleft{}rain\textquotedblright{} (Monier-Williams): the rain of the season.
\end{cousinweb}

\begin{grammarlens}[title={What you've built}]
One more word for getting around.
\end{grammarlens}

\subsection*{Your turn}
\begin{itemize}
\raggedright
  \item \emph{Say it:} \emph{kālavarṣaṁ}
  \item \emph{Say it:} \emph{kālavarṣaṁ}, once more
  \item \emph{Say it:} \emph{kālavarṣaṁ ā varṣaṁ nēratte etti.} — The southwest monsoon came early that year.
  \item \emph{Recall:} say the Malayalam for rain clouds, then the Malayalam for the monsoon, then say \emph{kālavarṣaṁ} again
\end{itemize}

\begin{tcolorbox}[breakable,colback=teal!4,colframe=teal!35!black,title={Before you move on}]
What does \ml{കാലവർഷം} mean? (\textquotedblleft{}The southwest monsoon\textquotedblright{}.) Say it once more.
\end{tcolorbox}
\section[\texorpdfstring{\ml{തുലാവർഷം} (tulāvarṣaṁ)}{തുലാവർഷം (tulāvarṣaṁ)}]{\ml{തുലാവർഷം} (tulāvarṣaṁ) — the northeast monsoon}

\label{lesson:ML-C507-tulavarsam}



\begin{quote}
Before the new one: say the Malayalam for the monsoon, then the Malayalam for the southwest monsoon.
\end{quote}

\subsection*{You'll want to know: \ml{തുലാവർഷം}}
\textbf{\ml{തുലാവർഷം}} — \emph{tulāvarṣaṁ} — \textquotedblleft{}the northeast monsoon\textquotedblright{}.

\begin{quote}
\ml{തുലാവർഷത്തിൽ} \ml{വൈകുന്നേരം} \ml{ഇടിമിന്നലോടെ} \ml{മഴ} \ml{പെയ്യും}.
\end{quote}

\begin{quote}
\emph{tulāvarṣattil vaikunnēraṁ iṭiminnalōṭe maḻa peyyuṁ.}
\end{quote}

\begin{quote}
In the northeast monsoon it rains in the evenings with thunder.
\end{quote}

It comes in the month of \ml{തുലാം}, October to November.

\begin{cousinweb}
\textbf{\ml{തുലാവർഷം}} (\emph{tulāvarṣaṁ}) is a Malayalam compound built from \ml{തുലാം}, the month named after Sanskrit \emph{tulā}, the scales, and \emph{varṣa}, rain (Monier-Williams).
\end{cousinweb}

\begin{grammarlens}[title={What you've built}]
One more word for getting around.
\end{grammarlens}

\subsection*{Your turn}
\begin{itemize}
\raggedright
  \item \emph{Say it:} \emph{tulāvarṣaṁ}
  \item \emph{Say it:} \emph{tulāvarṣaṁ}, once more
  \item \emph{Say it:} \emph{tulāvarṣattil vaikunnēraṁ iṭiminnalōṭe maḻa peyyuṁ.} — In the northeast monsoon it rains in the evenings with thunder.
  \item \emph{Recall:} say the Malayalam for the monsoon, then the Malayalam for the southwest monsoon, then say \emph{tulāvarṣaṁ} again
\end{itemize}

\begin{tcolorbox}[breakable,colback=teal!4,colframe=teal!35!black,title={Before you move on}]
What does \ml{തുലാവർഷം} mean? (\textquotedblleft{}The northeast monsoon\textquotedblright{}.) Say it once more.
\end{tcolorbox}
\section[\texorpdfstring{\ml{തോരുക} (tōruka)}{തോരുക (tōruka)}]{\ml{തോരുക} (tōruka) — to stop raining; to dry}

\label{lesson:ML-C507-toruka}



\begin{quote}
Before the new one: say the Malayalam for the southwest monsoon, then the Malayalam for the northeast monsoon.
\end{quote}

\subsection*{You'll want to know: \ml{തോരുക}}
\textbf{\ml{തോരുക}} — \emph{tōruka} — \textquotedblleft{}to stop raining; to dry\textquotedblright{}.

\begin{quote}
\ml{ഉച്ചയോടെ} \ml{മഴ} \ml{തോർന്നു}.
\end{quote}

\begin{quote}
\emph{uccayōṭe maḻa tōrnnu.}
\end{quote}

\begin{quote}
By noon the rain had stopped.
\end{quote}

Past: \ml{തോർന്നു}, \emph{tōrnnu}.

For rain, tears and wet things: \ml{കണ്ണീർ} \ml{തോർന്നു}, the tears dried.

\begin{cousinweb}
\textbf{\ml{തോരുക}} (\emph{tōruka}) is native Dravidian, from a root that means to dry or to wipe off moisture (DEDR 3351). DEDR also lists under it a word for fair weather when the monsoon rain stops.
\end{cousinweb}

\begin{grammarlens}[title={What you've built}]
One more word for getting around.
\end{grammarlens}

\subsection*{Your turn}
\begin{itemize}
\raggedright
  \item \emph{Say it:} \emph{tōruka}
  \item \emph{Say it:} \emph{tōruka}, once more
  \item \emph{Say it:} \emph{uccayōṭe maḻa tōrnnu.} — By noon the rain had stopped.
  \item \emph{Recall:} say the Malayalam for the southwest monsoon, then the Malayalam for the northeast monsoon, then say \emph{tōruka} again
\end{itemize}

\begin{tcolorbox}[breakable,colback=teal!4,colframe=teal!35!black,title={Before you move on}]
What does \ml{തോരുക} mean? (\textquotedblleft{}To stop raining; to dry\textquotedblright{}.) Say it once more.
\end{tcolorbox}
\section[\texorpdfstring{\ml{ചാറ്റൽമഴ} (cāṟṟalmaḻa)}{ചാറ്റൽമഴ (cāṟṟalmaḻa)}]{\ml{ചാറ്റൽമഴ} (cāṟṟalmaḻa) — a light drizzle}

\label{lesson:ML-C507-carralmala}



\begin{quote}
Before the new one: say the Malayalam for the northeast monsoon, then the Malayalam for to stop raining.
\end{quote}

\subsection*{You'll want to know: \ml{ചാറ്റൽമഴ}}
\textbf{\ml{ചാറ്റൽമഴ}} — \emph{cāṟṟalmaḻa} — \textquotedblleft{}a light drizzle\textquotedblright{}.

\begin{quote}
\ml{രാവിലെ} \ml{മുതൽ} \ml{ചാറ്റൽമഴ} \ml{ഉണ്ട്}.
\end{quote}

\begin{quote}
\emph{rāvile mutal cāṟṟalmaḻa uṇṭŭ.}
\end{quote}

\begin{quote}
There has been a light drizzle since morning.
\end{quote}

\ml{ചാറ്റൽ}, drizzling, plus \ml{മഴ}, rain.

\begin{cousinweb}
\textbf{\ml{ചാറ്റൽമഴ}} (\emph{cāṟṟalmaḻa}) is a native Dravidian compound of \ml{ചാറ്റൽ}, drizzling (DEDR 2482), and \ml{മഴ}, rain (DEDR 4753).
\end{cousinweb}

\begin{grammarlens}[title={What you've built}]
One more word for getting around.
\end{grammarlens}

\subsection*{Your turn}
\begin{itemize}
\raggedright
  \item \emph{Say it:} \emph{cāṟṟalmaḻa}
  \item \emph{Say it:} \emph{cāṟṟalmaḻa}, once more
  \item \emph{Say it:} \emph{rāvile mutal cāṟṟalmaḻa uṇṭŭ.} — There has been a light drizzle since morning.
  \item \emph{Recall:} say the Malayalam for the northeast monsoon, then the Malayalam for to stop raining, then say \emph{cāṟṟalmaḻa} again
\end{itemize}

\begin{tcolorbox}[breakable,colback=teal!4,colframe=teal!35!black,title={Before you move on}]
What does \ml{ചാറ്റൽമഴ} mean? (\textquotedblleft{}A light drizzle\textquotedblright{}.) Say it once more.
\end{tcolorbox}
